\section{Introduction}
\label{sec:intro}

FRB/US is the Federal Reserve Board's large-scale quarterly macroeconometric model of the United States. Introduced in 1996 as the successor to the MPS model, it is one of the principal models that Board staff consult for forecasting and for the analysis of monetary and fiscal policy issues \citep{brayton1996guide, brayton2014note}. Unusually among central-bank workhorse models, FRB/US is public: the Board publishes the complete equation system with estimated coefficients (\texttt{model.xml}), the quarterly historical-plus-projection database (\texttt{LONGBASE.TXT}), extensive technical documentation, and reference simulation platforms in EViews and, since 2021, in Python (the \texttt{pyfrbus} package), all in the public domain \citep{fedfrbusabout, fedfrbuspython}.

This paper documents \texttt{frbus}, an open, from-scratch reimplementation of the FRB/US simulation platform in modern Python, developed as part of PolicyEngine Macro --- PolicyEngine's programme of open macroeconomic modelling infrastructure that complements its tax--benefit microsimulation models. The motivation is practical. The Fed's own \texttt{pyfrbus} reference package is an invaluable specification of the model's semantics, but it is pinned to an ageing scientific-Python stack (\texttt{sympy} 1.3, \texttt{scikit-umfpack}, pre-2.0 \texttt{pandas}) and does not install cleanly on current interpreters without patching. A dependable, continuously tested implementation on current \texttt{numpy}/\texttt{scipy}/\texttt{sympy}/\texttt{pandas} lowers the barrier to using FRB/US in research and in downstream policy tooling, where microsimulation-based revenue and distributional estimates can be combined with model-consistent macroeconomic feedbacks.

The reimplementation is deliberately conservative: it treats the Fed's published \texttt{model.xml} and LONGBASE data as the single source of truth and follows \texttt{pyfrbus} semantics exactly, while making independent numerical choices where these are equivalence-preserving (a single-block Newton solve of the full simultaneous system per period, rather than the vendor's block decomposition). Fidelity is then demonstrated rather than asserted, in two layers. First, a \emph{tracking invariant}: after computing add factors, solving the model over 2026Q1--2030Q4 reproduces the LONGBASE baseline for all 284 endogenous variables to $5.6\times10^{-17}$ --- machine precision. Second, \emph{cross-validation}: the identical monetary-policy experiment run in this implementation and in the Fed's \texttt{pyfrbus} (in a separate, patched environment) agrees across all 284 endogenous variables and all 20 quarters to a maximum absolute difference of $6.0\times10^{-9}$, with the headline variables --- real GDP, prices, unemployment, the funds rate --- agreeing to $10^{-10}$ or better. Both tests run as hard gates in continuous integration.

Beyond replication of the reference implementation, the model's simulation properties line up with what the Fed has published about FRB/US itself. A 100-basis-point shock to the inertial Taylor rule lowers real GDP by 0.55 per cent at its trough seven quarters out, raises the unemployment rate by up to 0.26 percentage point, and lowers core PCE inflation by up to 0.034 percentage point --- signs, magnitudes and timing consistent with the funds-rate-shock responses under VAR expectations reported by \citet{brayton2014note}. A simple government-purchases experiment yields a first-year spending multiplier of about 0.7 with monetary policy responding through the inertial Taylor rule, in line with the model's documented fiscal properties.

The remainder of the paper proceeds as follows. Section~\ref{sec:model} summarises the history and structure of FRB/US, including the polynomial-adjustment-cost (PAC) framework and the expectations options, with representative equations. Section~\ref{sec:implementation} describes the implementation. Section~\ref{sec:validation} presents the validation evidence. Section~\ref{sec:results} presents simulation results and compares them with published FRB/US properties. Section~\ref{sec:limitations} discusses limitations.
